\documentclass[../../main.tex]{subfiles}
\begin{document}

\begin{flushright}
\footnotesize\textit{【自動文字起こし・要確認】}
\end{flushright}

{\bf [解]}

$\alpha, \beta$は独立に$1,\dots,6$の一様分布だから，その期待値は
\begin{align}
 E(\alpha) = E(\beta) = \dfrac{1}{6}\sum_{k=1}^{6}k = \dfrac{7}{2} \label{eq:1}
\end{align}
である．あとで使うため二乗の期待値も計算すると
\begin{align}
 E(\alpha^2)=E(\beta^2)=\dfrac{1}{6}\displaystyle\sum_{k=1}^6k^2=\dfrac{91}{6} \label{eq:2}
\end{align}
である．


\paragraph{(1)}

解と係数の関係より，$s,t$は$\alpha,\beta$を用いて
\begin{align}
 s=-(\alpha+\beta),\qquad t=\alpha\beta \label{eq:3}
\end{align}
と表せる．
$\alpha,\beta$は独立だからこれらの期待値は\cref{eq:1}より
\begin{align}
 \begin{cases}
 E(s)=E(-\alpha-\beta) = -E(\alpha)-E(\beta)=-7 \\
 E(t)=E(\alpha\beta)   = E(\alpha)E(\beta) =\dfrac{49}{4}  
 \end{cases} \label{eq:4}
\end{align}
である．

\paragraph{(2)}

題意より
\begin{align}
 f(x)^2=(x^2+sx+t)^2=x^4+2sx^3+(s^2+2t)x^2+2stx+t^2
\end{align}
だから係数を比較すると
\begin{align}
a=2s,\qquad b=s^2+2t,\qquad c=2st,\qquad d=t^2
\end{align}
である．ここに\cref{eq:3}を代入して（後々のため$a$のみそのまま残している．）
\begin{align}
 \begin{cases}
  a = 2s \\
  b = (\alpha+\beta)^2+2\alpha\beta=\alpha^2+4\alpha\beta+\beta^2 \\
  c = 2st=-2(\alpha+\beta)\alpha\beta=-2(\alpha^2\beta+\alpha\beta^2) \\
  d = t^2=\alpha^2\beta^2
 \end{cases} \label{eq:5}
\end{align}
である．

\subparagraph{$a$について}

\cref{eq:4,eq:5}より
\begin{align}
E(a)=2E(s)=-14
\end{align}
である．

\subparagraph{$b$について}

\cref{eq:5,eq:1}と$\alpha,\beta$の独立性より
\begin{align}
 E(b)
  &=E(\alpha^2)+4E(\alpha)E(\beta)+E(\beta^2) \\
  &=2\cdot\dfrac{91}{6}+4\cdot\dfrac{49}{4} \\
  &=\dfrac{91}{3}+49=\dfrac{238}{3}
\end{align}
である．

\subparagraph{$c$について}

\cref{eq:1,eq:5}と$\alpha,\beta$の独立性より
\begin{align}
 E(c)
  &=-2\bigl\{E(\alpha^2)E(\beta)+E(\alpha)E(\beta^2)\bigr\} \\
  &=-2\cdot2\cdot\dfrac{91}{6}\cdot\dfrac{7}{2} \\
  &=-\dfrac{637}{3}
\end{align}
である．

\subparagraph{$d$について}

\cref{eq:1,eq:5}と$\alpha,\beta$の独立性より
\begin{align}
 E(d)
  &=E(\alpha^2)E(\beta^2) \\
  &=\left(\dfrac{91}{6}\right)^2 \\
  &=\dfrac{8281}{36}
\end{align}
である．
\casesend

以上をまとめて
\begin{align}
 \begin{cases}
  E(a)=-14 \\
  E(b)=\dfrac{238}{3} \\
  E(c)=-\dfrac{637}{3} \\
  E(d)=\dfrac{8281}{36}
 \end{cases}
\end{align}
である．

\newpage

\end{document}
